\section{Limits and Failure Modes}

The implicit EM framework is not universal. It applies under specific conditions, and when those conditions fail, so does the analysis. This section delineates the boundaries: when implicit EM does not arise, what pathologies can occur even when it does, and what phenomena lie entirely outside its scope. A clear account of limits strengthens rather than weakens the contribution---it distinguishes precise claims from overreach.

\subsection{When Implicit EM Does Not Arise}

The identity $\partial L / \partial d_j = -r_j$ depends on the log-sum-exp structure. Remove normalization, and the identity does not hold.

Consider multi-label classification with independent sigmoids. Each class $j$ has its own output $\sigma(z_j)$, and the loss decomposes as a sum of independent binary cross-entropies. There is no partition function, no softmax, no competition across classes. The gradient with respect to $z_j$ depends only on the target for class $j$, not on the outputs of other classes.

In this setting, responsibilities do not exist. No quantity sums to one across classes; no soft assignment distributes an input among alternatives. Each output channel operates in isolation. A point can be equally close to all prototypes or far from all of them, and the gradients do not redistribute---they simply reflect independent errors.

This is not a failure of the architecture but an absence of the required objective structure. Implicit EM arises from competition, and competition arises from normalization. Systems with independent outputs may learn useful representations, but they do not perform mixture inference and do not exhibit responsibility-weighted dynamics. The explanatory scope of the framework ends where normalization ends.

\subsection{Scale and Collapse}

A full Gaussian mixture model includes a log-determinant term in the likelihood---a penalty on the volume of each component's covariance. This term prevents collapse: without it, a component could shrink its covariance to zero, placing infinite density on a single point and achieving unbounded likelihood. The log-determinant diverges as the covariance collapses, counterbalancing the density gain.

Most neural objectives omit this term. Cross-entropy and attention softmax operate on distances or scores without explicit volume penalties. The implicit EM dynamics still hold---gradients are still responsibility-weighted---but nothing prevents the learned metric from degenerating. A network can learn to map all inputs to nearby points, collapsing the distance structure and trivializing the responsibilities. (Section 3.4 shows that for linear-layer distances the exact volume term is available in network primitives at the cost of one log-determinant; the analysis here concerns the common case where it is omitted.)

In practice, collapse is often avoided through implicit mechanisms: weight decay regularizes the scale of projections; layer normalization constrains activation magnitudes; architectural choices like residual connections preserve signal diversity. These interventions work, but they are not derived from the objective---they are heuristics that happen to stabilize geometry.

The implicit EM framework clarifies why collapse is a risk. When components are updated in proportion to their responsibilities, a component that captures slightly more mass receives stronger gradients, captures more mass still, and can dominate entirely. This positive feedback is inherent to EM dynamics and is classically controlled by the volume term. Neural networks remove that control and rely on other mechanisms to fill the gap. The framework does not solve this problem; it explains why the problem exists.

\subsection{Supervision Constraints}

In the unsupervised regime, responsibilities are entirely latent---the data alone determines which components claim which inputs. Supervision changes this. Labels declare which component \emph{should} take responsibility, overriding what the geometry would otherwise dictate.

This constraint is powerful but rigid. Cross-entropy training forces the correct class toward responsibility 1, regardless of whether the input lies near that class's prototype or far from all prototypes. An input equidistant from all class boundaries still receives a hard label; the model must assign it somewhere. The soft, graded structure of responsibilities persists among incorrect classes, but the correct class is clamped.

A consequence is the closed-world assumption. Softmax normalization guarantees that responsibilities sum to one---some class must take full responsibility for every input. There is no ``none of the above,'' no option for the model to reject an input as foreign to all known categories. An out-of-distribution input, no matter how anomalous, will be assigned to whichever class has the lowest distance, and the model's confidence may be arbitrarily high.

This is a limitation of the objective, not the mechanism. The implicit EM framework describes how responsibilities arise and how they weight learning. It does not claim that the resulting assignments are semantically correct or that the objective captures all desirable behaviors. Cross-entropy with softmax enforces assignment; if non-assignment is needed, the objective must change. Open-set recognition, out-of-distribution detection, and selective prediction require objectives that break the closed-world assumption---either by introducing an explicit reject option or by abandoning normalization entirely. The implicit EM analysis applies to whatever objective is chosen; it does not choose the objective.

\subsection{What This Framework Does Not Explain}

The implicit EM framework explains one phenomenon: the emergence of responsibility-weighted learning dynamics from distance-based objectives. It does not explain everything neural networks do.

In-context, inference-time Bayesian computation is not addressed. A trained transformer with frozen weights can track a posterior over latent task variables as a sequence unfolds \citep{agarwal2025geometry,agarwal2025scaling}; that is a property of the learned function, executed in activation space at inference time. Our identity concerns gradient updates to parameters during training. The two are connected---training under cross-entropy drives the network toward the Bayes posterior predictive, and the training dynamics build the geometry the frozen network uses---but the framework derives the per-step structure of learning, not the algorithm the learned function implements in context. The forward-pass log-sum-exp gradient of Section 4.2 describes a single attention read-out, not sequential belief updating across a context.

Generalization---why networks perform well on unseen data---is not addressed. The framework describes the dynamics of training, not the inductive biases that enable generalization beyond the training distribution. Scaling laws---the predictable relationship between model size, data, and performance---lie outside the analysis entirely. Nothing in the gradient identity suggests how performance should scale with parameters or compute.

Long-horizon reasoning, planning, and sequential decision-making involve temporal structure that the framework does not capture. Implicit EM describes how a single input is softly assigned to components and how those components update. It does not describe how representations are composed over time, how goals propagate backward through action sequences, or how models learn to search.

Emergent capabilities---the sudden appearance of qualitatively new behaviors at scale---remain unexplained. The framework offers no account of why certain abilities appear discontinuously or why they require specific thresholds of model size. If emergent capabilities arise from implicit EM dynamics, the connection is not obvious; if they arise from other mechanisms, the framework is silent.

These are not failures of the analysis but boundaries of its scope. The contribution is to identify and derive one mechanism precisely, not to explain all of deep learning. Clarity about what is claimed protects against overinterpretation---and leaves room for complementary explanations of phenomena that implicit EM does not reach.
